\documentclass[11pt]{article}
\usepackage[a4paper,margin=1in]{geometry}
\usepackage{amsmath,amssymb,amsthm,graphicx,microtype,xcolor}
\usepackage[colorlinks=true,linkcolor=blue!50!black,urlcolor=blue!60!black,citecolor=blue!50!black]{hyperref}
\newtheorem{theorem}{Theorem}
\newtheorem{lemma}{Lemma}
\newcommand{\R}{\mathbb R}
\newcommand{\C}{\mathbb C}
\newcommand{\Sph}{\mathbb S^1}
\newcommand{\PhiS}{\Phi^{\mathcal S}}
\newcommand{\PhiA}{\Phi^{\mathcal A}}
\title{A vortex obstruction to openness of the wave-function--density map\\
\large A counterexample to the unqualified question in arXiv:1907.02024}
\author{Autonomous mathematical discovery run \texttt{fa\_banach\_001}, lane 5}
\date{12 August 2026}

\begin{document}
\maketitle

\begin{center}
\fbox{\parbox{0.91\linewidth}{\textbf{Status: candidate counterexample, likely valid.}
The construction is exact and analytic.  It settles the question as stated without an $N\geq2$ restriction, already for one spinless particle.}}
\end{center}

\section{Source question}

In \emph{From wave-functions to single electron densities}, Ugo Bindini and Luigi De Pascale consider the maps from symmetric and antisymmetric normalized $H^1$ wave functions to their one-particle densities \cite[p.~2]{source}.  They ask:
\begin{quote}
Suppose that $\sqrt{\rho_k}\to\sqrt\rho$ in $H^1$, and take $\psi$ such that $\rho=\rho[\psi]$.  Can one find $\psi_k$ such that $\rho_k=\rho[\psi_k]$ and $\psi_k\to\psi$ in $H^1$?  In other words, are $\PhiA$ and $\PhiS$ open?
\end{quote}
Their Theorems 1.1--1.2 give positive answers for nonnegative and real-valued bosonic wave functions.  The 2026 journal version, retitled \emph{From Wave-Functions to Single Boson Densities}, still explicitly describes the general question as open \cite{journal}.  Figure~\ref{fig:source} shows the source formulation and its partial results.

\begin{figure}[ht]
\centering
\includegraphics[width=0.94\linewidth]{figures/open_problem_crop.png}
\caption{The open question and partial positive results on p.~2 of the source paper.}
\label{fig:source}
\end{figure}

\section{Proof intuition}

For one spinless particle, the density is simply $\rho[\psi]=|\psi|^2$.  A complex wave function can carry phase topology invisible to its density.  Take a compactly supported wave function whose phase winds once around the $x_1x_2$-axis.  Its modulus vanishes on that axis.  Fill the zero with an arbitrarily small positive core; the new square-root density converges in $H^1$, but a wave function having that positive modulus has an $H^1$ circle-valued phase across every transverse disk.  Such a disk phase has degree zero on interior circles.  It therefore cannot converge in $H^1$ to the degree-one vortex phase.

The point is specific and sharp: the source's positive theorem covers real phases, while the obstruction is genuinely complex.  It is also invisible at the density level.

\section{A Sobolev phase lemma}

\begin{lemma}[No $H^1$ vortex without a zero]\label{lem:lifting}
Let $D\subset\R^2$ be a disk and $v\in H^1(D;\Sph)$.  Then $v=e^{i\varphi}$ for some $\varphi\in H^1(D;\R)$.  Consequently, for almost every concentric circle $\partial D_r\subset D$, the restriction of $v$ has degree zero.
\end{lemma}

\begin{proof}
Set $j=-i\overline v\,\nabla v\in L^2(D;\R^2)$.  Differentiating $|v|^2=1$ shows that each weak derivative of $v$ is tangent to $\Sph$, so $\partial_kv=i j_kv$ almost everywhere.  The distributional curl of $j$ is zero: after cancellation of mixed derivatives, the remaining expression is
\[
 -i(\partial_1\overline v\,\partial_2v-\partial_2\overline v\,\partial_1v)=0.
\]
The Sobolev Poincar\'e lemma on the simply connected disk gives $\varphi\in H^1(D)$ with $\nabla\varphi=j$.  Hence $\nabla(e^{-i\varphi}v)=0$, and changing $\varphi$ by a constant yields $v=e^{i\varphi}$.

For almost every $r$, the circular trace $\varphi(r,\cdot)$ belongs to $H^1(\Sph)$ and is periodic.  Therefore
\[
 \deg(v|_{\partial D_r})
 =\frac{1}{2\pi i}\int_0^{2\pi}\overline v\,\partial_\theta v\,d\theta
 =\frac{1}{2\pi}\int_0^{2\pi}\partial_\theta\varphi\,d\theta=0.
\]
\end{proof}

\section{Counterexample}

\begin{theorem}\label{thm:counterexample}
In the spinless one-particle sector, for every spatial dimension $d\geq2$, neither the symmetric nor the antisymmetric wave-function--density map is open in the topology used in \cite{source}.  More precisely, there are a normalized $\psi\in H^1(\R^d;\C)$ and probability densities $\rho_n$ such that
\[
 \sqrt{\rho_n}\longrightarrow\sqrt{\rho[\psi]}\quad\text{in }H^1(\R^d),
\]
but no choice of $\psi_n\in H^1(\R^d;\C)$ satisfying $\rho[\psi_n]=\rho_n$ can converge to $\psi$ in $H^1$.
\end{theorem}

\begin{proof}
Write $x=(z,y)$ with $z=(x_1,x_2)\in\R^2$ and $y\in\R^{d-2}$ (omit $y$ when $d=2$).  Choose nonnegative smooth compactly supported cutoffs $\chi(z)$ and $\zeta(y)$ which equal one on $|z|\leq3$ and on a nonempty cube $Q\subset\R^{d-2}$, respectively, with $\chi$ radial.  For $d=2$, set $\zeta\equiv1$.  Let
\[
 \psi(z,y)=C_0\,(x_1+i x_2)\chi(z)\zeta(y),
\]
where $C_0>0$ normalizes the $L^2$ norm to one.  Its modulus is
\[
 a(z,y):=|\psi(z,y)|=C_0|z|\chi(z)\zeta(y).
\]
For $\varepsilon>0$, define
\[
 a_\varepsilon(z,y)=C_\varepsilon
 \sqrt{|z|^2+\varepsilon^2}\,\chi(z)\zeta(y),
 \qquad \rho_\varepsilon=a_\varepsilon^2,
\]
where $C_\varepsilon$ makes $\|a_\varepsilon\|_2=1$.  Plainly $C_\varepsilon\to C_0$.  Moreover,
\[
 \nabla_z\sqrt{|z|^2+\varepsilon^2}
 =\frac{z}{\sqrt{|z|^2+\varepsilon^2}}\longrightarrow\frac z{|z|}
 \quad\text{a.e.},
\]
and these gradients are bounded by one.  Applying dominated convergence also to the cutoff terms gives
\begin{equation}\label{eq:H1root}
 a_\varepsilon\longrightarrow a=|\psi|\quad\text{strongly in }H^1(\R^d).
\end{equation}
Thus $\rho_\varepsilon$ belongs to the source's admissible density class and has the required convergence.

Take any sequence $\varepsilon_n\downarrow0$ and suppose, for contradiction, that $\psi_n\to\psi$ in $H^1$ while $|\psi_n|=a_{\varepsilon_n}$ almost everywhere.  On the cylinder
\[
 \Omega=\{\,|z|<3\,\}\times Q
\]
the modulus $a_{\varepsilon_n}$ is strictly positive.  Hence
\[
 u_n:=\frac{\psi_n}{a_{\varepsilon_n}}\in H^1(\Omega;\Sph).
\]
On the annular cylinder
\[
 \Omega'=\{\,1<|z|<2\,\}\times Q,
\]
the functions $a_{\varepsilon_n}$ and their reciprocals converge in $W^{1,\infty}$ to $C_0|z|$ and $(C_0|z|)^{-1}$.  It follows from $\psi_n\to\psi$ that
\begin{equation}\label{eq:phaseconv}
 u_n\longrightarrow u(z,y):=\frac{x_1+i x_2}{|z|}
 \quad\text{strongly in }H^1(\Omega').
\end{equation}

Pass to a subsequence for which the squared $H^1(\Omega')$ errors in \eqref{eq:phaseconv} are summable.  Fubini's theorem first in $y$ and then in the polar radius produces one $y_0\in Q$ (unneeded for $d=2$) and one $r_0\in(1,2)$ such that
\[
 u_n(r_0,\cdot,y_0)\longrightarrow e^{i\theta}
 \quad\text{strongly in }H^1(\Sph),
\]
and such that every disk slice $u_n(\cdot,y_0)$ lies in $H^1(\{|z|<3\};\Sph)$.  Lemma~\ref{lem:lifting} gives
\[
 \deg u_n(r_0,\cdot,y_0)=0
 \quad\text{for every }n.
\]
But degree is continuous under strong $H^1(\Sph)$ convergence: the formula
\[
 \deg w=\frac{1}{2\pi i}\int_{0}^{2\pi}\overline w\,w'\,d\theta
\]
passes to the limit because $H^1(\Sph)$ embeds into $C^0(\Sph)$.  The limit $e^{i\theta}$ has degree one, a contradiction.

For $N=1$ the permutation group is trivial, so the same $\psi$ belongs to both the symmetric and antisymmetric classes.  Finally, failure of the stated sequential lifting property implies non-openness: if either map were open at $\psi$, a diagonal choice from preimages of the approaching $\rho_{\varepsilon_n}$ would produce the forbidden convergent lifting sequence.
\end{proof}

\section{Verification, novelty, and scope}

Every step above is qualitative and exact; no numerical verification is needed.  The two points most worth checking are the $H^1$ lifting lemma and the Fubini selection of a common two-dimensional slice and circle.  The latter is standard: summability of the annular $H^1$ errors makes the sliced errors tend to zero for almost every $(y,r)$, while the trace and lifting conclusions fail only on countable unions of null sets.

A bounded novelty search on 12 August 2026 covered the run's registry, solution, attempt, and proof-gap indexes; arXiv and current web searches for the exact title and question; and combinations of ``Sobolev modulus map,'' ``$H^1$ vortex degree,'' ``Lieb Question 2,'' and ``counterexample.''  No prior counterexample was located.  The May 2026 journal article still calls the general question open and says that its vectorial/phase step is incomplete \cite{journal}.  This supports, but cannot guarantee, novelty.

\paragraph{Scope limitation.}
The theorem disproves the question as formulated without excluding one-particle systems.  It does \emph{not} settle a reformulation imposing $N\geq2$.  In a many-particle disintegration, the normalized conditional wave function takes values in an infinite-dimensional unit sphere, so the scalar $\Sph$ degree obstruction used here need not persist.  Human review should first confirm that $N=1$ lies within the intended quantifiers of Lieb's and the source paper's question.  If an implicit many-particle convention excludes it, this remains a sharp one-particle obstruction rather than a full answer to that restricted variant.

\sloppy
\begin{thebibliography}{9}
\bibitem{source}
U. Bindini and L. De Pascale,
\emph{From wave-functions to single electron densities},
arXiv:1907.02024, 2019.
\url{https://arxiv.org/abs/1907.02024}.

\bibitem{journal}
U. Bindini and L. De Pascale,
\emph{From Wave-Functions to Single Boson Densities},
Archive for Rational Mechanics and Analysis 250 (2026), Article 44.
\url{https://doi.org/10.1007/s00205-026-02207-2}.

\bibitem{lieb}
E. H. Lieb,
\emph{Density functionals for Coulomb systems},
International Journal of Quantum Chemistry 24 (1983), 243--277.
\end{thebibliography}

\end{document}
